%% kapa: exam=1 round=23 year=2012 period=1 form=A subject=economics subject_ko=경제원론 seq=20 number=60
%% source: Q-Net 감정평가사 기출문제 (cst00309) artlSeq=5052578 file=raw/1008/5052578/2012년 제23회 감정평가사 제1차시험 기출문제(1교시 A형).pdf
%% conversion: kordoc 4.24.51 -> tools/gen1.py (auto); review=proofread-sonnet; flags=none
\begin{question}{60}
A국은 자본이 상대적으로 풍부하고 B국은 노동이 상대적으로 풍부하다. 양국 간의 상품이동이 완전히 자유로워지고 양 국가가 부분 특화하는 경우, 헥셔-올린(Heckscher-Ohlin)모형과 스톨퍼-새뮤엘슨(Stolper-Samuelson)정리에서의 결과와 부합하는 것을 모두 고른 것은?
\begin{framedtext}
ㄱ. 두 국가의 자본가격은 같아진다. \par
ㄴ. B국 자본가의 실질소득이 증가한다. \par
ㄷ. A국 노동자의 실질소득이 감소하는 반면, B국 노동자의 실질소득은 증가한다.
\end{framedtext}
\begin{choices}
  \item ㄱ
  \item ㄱ, ㄴ
  \item ㄱ, ㄷ
  \item ㄴ, ㄷ
  \item ㄱ, ㄴ, ㄷ
\end{choices}
\end{question}
\answer{3}
